%%%%%%%%%%%%%%%%%%%%%%%%%%%%%%%%%%%%%%%%%%%%%%%%%%%%%%%%%%%%%%%%%%%%%%
\documentclass[letterpaper]{article}
\usepackage{latexsym}
\usepackage{graphicx}
\usepackage{amsmath}

\begin{document}
\title{Lunar Mission Software: Launch Mission Software Development
Initial Writeup}
\author{Lou Rosas}
\maketitle
\section{History}
Lunar Mission Software is complex and entails many different phases
(states).  Attempts to tackle the entire software as a single 
development cycle is proven futile and disastrous.  The best 
approach is break the Software into ``Meta States" that handles a 
given phase of the Development--From Launch to Recovery.  The Software
is then integrated as a larger System.

\section{Abstract}
Using modern Software Development Techniques, develop Launch Mission 
Software for integrating into The Lunar Mission Software.  The 
Software is independent of the Launch Mission Software in the sense 
the ability to run in a standalone mode.

\section{Concept}
To gain understanding in the full development of the Launch Mission
Software.  Including Analysis and Design techniques.  Including 
gaining understanding in concurrency development.

\section{Intent}
As part of the Lunar Mission Software Development, develop the Launch
Mission Software.  The requesting System shall be ``input blind" and
the shall not be able to distinguish between software stimuli and 
actual launch hardware input.

\section{Stakeholders \& Interests}
\begin{itemize}
   \item Flight Director Wants to see a successful launch without
      incident
   \item Flight Engineers \& Technicians Wants accurate data for the
      components of the System responsible for monitoring,
   \item Mission Planners Want to ensure a successful mission and
      launch
   \item Capsule Crew/Astronauts Want a successful launch
\end{itemize}

\section{Typical Success Scenario}
\begin{enumerate}
   \item The Rocket with payload/capsule is loaded onto the Launching
      Mechanism/Platform and is moved to the Launching site (if not 
      already there)--it is not fueled...
   \item The Flight Director loads the the Rocket Profile into the
      system this includes the rocket data, the stage data, the 
      capsule data, launch platform data--This is the intended Rocket
      to Launch as related to the mission
   \item The System Loads in the Parameters per the current Launch
      \begin{enumerate}
         \item This starts the \textit{Initialization State}
      \end{enumerate}
   \item The System Begins monitoring the Components pertinent to the
      Launch (System Components like the Rocket, Launch Platform, 
      et. al.) The Components begin real time feeding data to the
      System
   \item If the System Reports errors in any of the System, the 
      Flight Director has the Specialists investigate, evaluate and as
      needed correct any errors
   \item Once the Flight Director has determined the System is Safe,
      the Flight Director determines the Countdown Time, enters it and
      starts the Countdown
      \begin{enumerate}
         \item \textit{This transitions to the Countdown State}
         \item The System Begins Counting Down to the Launch
      \end{enumerate}
   \item During the Countdown, the Rocket is loaded with fuel
      \begin{enumerate}
         \item Each Stage is fueled
         \item \textit{This transitions to the Fueling Substate within
            the Countdown State}
      \end{enumerate}
   \item Fueling continues until each Stage is full
      \begin{enumerate}
         \item The System continuously monitors the fueling.
         \item Fueling  and the countdown stops if any anomalies or
            errors are reported.
            \begin{enumerate}
               \item \textit{This transitions to the Stop State}
            \end{enumerate}
         \item  Fueling does not continue until the errors and
            anomalies are fixed
            \begin{enumerate}
               \item If the errors are severe enough to where it is
                  impossible to fix without aborting the launch, then
                  the Countdown is aborted
                  \begin{enumerate}
                     \item \textit{This transitions to the Abort
                        State}
                  \end{enumerate}
            \end{enumerate}
         \item Once the errors are fixed, Fueling is resumed and the
            countdown continues
            \begin{enumerate}
               \item \textit{Transitions to the Countdown State}
               \item \textit{Transitions to the Fueling Substate
                  within the Countdown State}
            \end{enumerate}
         \item \textit{System Transitions out of the Fueling Substate
            into the Fueled Substate once Stages are Fueled}
      \end{enumerate}
   \item During the countdown, if any of the components report an 
      anomaly or an error, the countdown is stopped
      \begin{enumerate}
         \item \textit{The System Transitions to the Stop State}
         \item Investigated and if needed, fixed
            \begin{enumerate}
               \item If the errors are such that it is impossible to 
                  fix without aborting the Countdown, then the 
                  Countdown is aborted
                  \begin{enumerate}
                     \item \textit{The System Transitions to the Abort
                        State}
                  \end{enumerate}
               \item If the errors are fixed, the countdown continues
                  \begin{enumerate}
                     \item \textit{The System Transitions back to the
                        Countdown State and appropriate Substate}
                  \end{enumerate}
            \end{enumerate}
      \end{enumerate}
   \item Once the Countdown reaches 0, the Flight Director makes the 
      choice to ignite the rocket or abort the Launch.
      \begin{enumerate}
         \item \textit{Transitions to Ignition State if ignited}
         \item \textit{Transitions to Abort State if aborted}
      \end{enumerate}
   \item The Rocket continuously monitors the Thrust and Fuel Delivery
      System, reporting anomalies and or errors.
      \begin{enumerate}
         \item If there are errors, the Stakeholders have the option to
            Abort
         \item \textit{System Transitions to the Abort State}
      \end{enumerate}
\end{enumerate}
\end{document}
%%%%%%%%%%%%%%%%%%%%%%%%%%%%%%%%%%%%%%%%%%%%%%%%%%%%%%%%%%%%%%%%%%%%%%
